% Generated draft: requires mathematical and source review.
\documentclass[11pt]{article}
\usepackage[margin=25mm]{geometry}
\usepackage{fontspec}
\setmainfont{Latin Modern Roman}
\usepackage{amsmath,amssymb,mathtools,unicode-math}
\setmathfont{Latin Modern Math}
\usepackage{enumitem,needspace,xurl,hyperref}
\setlength{\parindent}{0pt}
\setlength{\parskip}{5pt}
\setlength{\emergencystretch}{3em}
\newcommand{\N}{\mathbb{N}}
\newcommand{\Z}{\mathbb{Z}}
\newcommand{\Q}{\mathbb{Q}}
\newcommand{\R}{\mathbb{R}}
\newcommand{\C}{\mathbb{C}}
\newcommand{\F}{\mathbb{F}}
\newcommand{\A}{\mathbb{A}}
\newcommand{\PP}{\mathbb{P}}
\newcommand{\E}{\mathbb{E}}
\newcommand{\eps}{\varepsilon}
\newcommand{\dd}{\,\mathrm{d}}
\begin{document}
\section*{UnsolvedMath ID: 30001010}
\textbf{Source problem code:} catalogue entry
\subsection{Definitions and mathematical statement}
Let $X$ be compact Hausdorff and let $E\to X$ be a complex vector bundle of constant positive rank $n$. Its continuous sections form a Hilbert $C(X)$-module, with pointwise left action, and the associated Cuntz--Pimsner algebra $O_E$ is the source's locally trivial $C(X)$-algebra with fiber $O_n$. The Cuntz algebra $O_n$ is generated by isometries $s_1,\ldots,s_n$ with $s_i^*s_j=\delta_{ij}1$ and $\sum_i s_is_i^*=1$; for $n=1$ this convention gives $C(S^1)$.

For fixed $X,n$, consider the association
\[
[E]\longmapsto[O_E].
\]
The question asks which invariants of the bundle can be recovered from the field's isomorphism class. In the natural base-preserving interpretation, characterize the equivalence relation
\[
E\sim F\iff O_E\cong O_F\text{ as }C(X)\text{-algebras},
\]
and determine which bundle invariants $q$ satisfy $E\sim F\Rightarrow q(E)=q(F)$. This states recovery of information; it does not demand that $O_E$ determine the full bundle isomorphism class.

The source does not specify whether isomorphism must preserve $C(X)$, the gauge action, or a bundle of fibers, or whether arbitrary algebra isomorphisms and base homeomorphisms are allowed. These categories can retain different information. Its supplied partial result says that for finite-dimensional $X$ the $K$-class $[E]\in K^0(X)$ determines $O_E$. That is a forward implication from bundle $K$-data to the field; it does not alone say that the field determines $[E]$. The general compact-Hausdorff scope is not replaced by a finite-dimensional or low-dimensional classification.
\subsection{Short English statement}
Determine which vector-bundle information can be recovered from its Cuntz--Pimsner field. The isomorphism category must be specified, and a result saying that a bundle $K$-class determines the field does not by itself prove the reverse recovery statement.
\subsection{Sources}
\begin{itemize}
\item Supplied report, with corrupted title markup normalized: \emph{$C^*$-Algebras} (2009), Oberwolfach Reports, DOI 10.4171/owr/2008/37, pp. 2095--2164; proposed year 2008. \url{https://doi.org/10.4171/owr/2008/37}.
\item Supplied Vasselli vector-bundle theorem reference: \url{https://doi.org/10.1016/j.jfa.2004.09.004}.
\item Local UnsolvedMath registry, permanent ID 30001010, catalogue entry; created 25 August 2026, updated 16 September 2026. The corrected\_verified note describes editorial cleanup. Literature checked at \texttt{2026-08-21}: the assessment is supplied dated curation, not independently verified. It labels the record partially solved, reports Vasselli's finite-dimensional $K$-class determination result and Dadarlat low-dimensional classifications, while retaining the broader invariant question. No separate Dadarlat bibliographic reference is supplied, so none is invented.
\end{itemize}
\end{document}
